\section{A spherical formula for the realized critical residual}
\label{sec:r13-spherical-critical-residual}

Let $T_0$ be a critical tangent supplied by MT on
$\mathbf R^N$, $N\ge2$.  Thus it is a symmetric,
middle-positive, separately closed real supercurrent of
bidimension $(2,2)$.  Put
\[
 \beta=\sum_i d'x_i\wedge d''x_i,\qquad
 \mu=\tfrac12T_0\wedge\beta^2,\qquad
 \mu(B_R)=\nu R^2.
 \tag{SR1}
\]
MT gives zero coefficient mass at the origin, conicity,
and radial two-vector support.  All norms and spheres in
this section are those of these fixed protected real
coordinates.  The conclusions are a polar reformulation
and corollaries of MT--LE--PT; they do not constitute a
new general relation-annihilation theorem.

\subsection{The angular matrix measure and its normalization}

Use the standard column-positive convention to associate
to a symmetric $(2,2)$ test form $\eta$ its coefficient
operator $C_\eta$ on $\bigwedge^2\mathbf R^N$.  In
particular
\[
                    C_{\beta^2/2}=I.
 \tag{SR2}
\]
For symmetric $(1,1)$ coefficient matrices $E,H$, this
convention is characterized by
\[
\begin{split}
C_{E\wedge H}(u\wedge v,s\wedge t)
={}&E(u,s)H(v,t)+H(u,s)E(v,t)\\
 &-E(u,t)H(v,s)-H(u,t)E(v,s).
\end{split}
 \tag{SR3}
\]
Let $\mathsf Q$ be the positive semidefinite matrix-valued
Radon measure representing $T_0$ in this convention:
\[
 \langle T_0,\eta\rangle
       =\int\operatorname{tr}(C_\eta\,d\mathsf Q),
                 \qquad \operatorname{tr}\mathsf Q=\mu.
 \tag{SR4}
\]
The trace is the full matrix contraction, including off-diagonal
entries.

Write $r=|x|$, $\theta=x/r$, and let
\[
 A_\theta:\mathbf R^N\longrightarrow\bigwedge^2\mathbf R^N,
                     \qquad A_\theta v=\theta\wedge v.
 \tag{SR5}
\]
It restricts to an isometry on $\theta^\perp$.  The radial
support assertion of MT says that $\mathsf Q$ is supported
as a matrix on $A_\theta(\theta^\perp)$ at $\mu$-almost
every nonzero point.

\begin{proposition}[Polar coefficient measure]
\label{prop:r13-angular-critical-matrix}
There is a unique finite positive semidefinite matrix-valued
measure $\mathsf M$ on $\mathbb S^{N-1}$ such that
\[
 \mathsf M\theta=0,\qquad
 \sigma:=\operatorname{tr}\mathsf M,\qquad
                   \sigma(\mathbb S^{N-1})=2\nu,
 \tag{SR6}
\]
and, in polar coordinates,
\[
 A_\theta^*\,d\mathsf Q\,A_\theta
                         =r\,dr\,d\mathsf M(\theta).
 \tag{SR7}
\]
The notation $\mathsf M\theta=0$ means that the matrix
density with respect to $\sigma$ kills $\theta$ almost
everywhere.  Equation SR7 determines $\mathsf Q$ completely,
since its matrix support is radial.
\end{proposition}

\begin{proof}
Put $d\mathsf B=A_\theta^*d\mathsf Q A_\theta$ on
$\mathbf R^N\setminus\{0\}$.  It is positive semidefinite,
kills $\theta$, and has trace $\mu$, by the radial support
and the isometry in SR5.  The coefficient form of MT's
conicity is
\[
             \mathsf Q(\lambda E)=\lambda^2\mathsf Q(E)
                              \quad(\lambda>0).
 \tag{SR8}
\]
Since $A_\theta$ has degree zero under dilation, the same
equation holds for $\mathsf B$.

For a Borel set $G\subset\mathbb S^{N-1}$ define
\[
 \mathsf M(G)=2\mathsf B
       \bigl(\{0<r\le1,\ \theta\in G\}\bigr).
 \tag{SR9}
\]
For $0<a<b$, SR8 gives
\[
 \mathsf B\bigl(\{a<r\le b,\ \theta\in G\}\bigr)
              =\frac{b^2-a^2}{2}\mathsf M(G).
\]
These identities on Borel rectangles prove SR7 by the
uniqueness of finite measures on compact radial annuli.
Positivity, tangency and SR6 follow from the corresponding
properties of $\mathsf B$ and $\mu(B_1)=\nu$.
The same rectangle identity proves uniqueness.
\end{proof}

\subsection{Radial cutoffs read a full angular tensor}

Let $E$ be a symmetric $(1,1)$ symbol whose coefficients
are homogeneous of degree zero and have a supplied genuine
compatible finite-PS ambient presentation on a conical
subdivision.  It is locally bounded.  Its restriction to the
sphere is denoted by the same letter.  Put
\[
 E_{rr}(\theta)=E(\theta,\theta),\qquad
            E_{\rm tan}=E|_{\theta^\perp\times\theta^\perp}.
 \tag{SR10}
\]
All face values use the canonical tangential PS convention
of PT.  Since $\mathsf M$ is tangent, the contraction
$E_{\rm tan}:d\mathsf M$ equals $E:d\mathsf M$.

\begin{theorem}[Angular formula]
\label{thm:r13-angular-critical-pairing}
Let $\rho(x)=\psi(|x|)$ be any compact smooth radial
function equal to one near the origin.  Then
\[
 \boxed{
 \langle T_0,E\wedge D\rho\rangle
       =\int_{\mathbb S^{N-1}}E_{\rm tan}:d\mathsf M
          -\int_{\mathbb S^{N-1}}E_{rr}\,d\sigma.}
 \tag{SR11}
\]
In particular the pairing is independent of the choice of
such a radial cutoff.  The angular pairings retain all
coefficient walls and all mixed entries.
\end{theorem}

\begin{proof}
On $r>0$ the matrix of $D\rho$ is
\[
       \psi''(r)\theta\theta^t
           +\frac{\psi'(r)}r(I-\theta\theta^t).
 \tag{SR12}
\]
For $u,v\perp\theta$, SR3 gives
\[
 C_{E\wedge D\rho}(\theta\wedge u,\theta\wedge v)
       =\psi''(r)E(u,v)
                +\frac{\psi'(r)}r E(\theta,\theta)(u,v).
 \tag{SR13}
\]
Here $(u,v)$ is the Euclidean inner product.  Radial
support removes all other coefficient blocks.  Combining
SR7 and SR13, and using the degree-zero homogeneity of
$E$, yields
\[
\begin{split}
\langle T_0,E\wedge D\rho\rangle
={}&\left(\int_0^\infty r\psi''(r)\,dr\right)
            \int_{\mathbb S^{N-1}}E_{\rm tan}:d\mathsf M\\
 &+\left(\int_0^\infty\psi'(r)\,dr\right)
            \int_{\mathbb S^{N-1}}E_{rr}\,d\sigma.
\end{split}
\]
All radial derivatives are supported in a compact annulus.
Their two integrals are $1$ and $-1$, respectively, because
$\psi=1$ near zero and vanishes near infinity.  This proves
SR11.

The use of finite-PS coefficients is legitimate by PT.
More explicitly, on a conical coefficient wall the normal
cocoefficients of $\mathsf Q$ vanish as measures.  Because
$\theta$ is tangent to that wall, SR5 transfers the same
annihilation to the corresponding normal entries of
$\mathsf M$.  Thus the displayed contractions do not depend
on ambient representatives of the face values.  Tangential
mass on the wall remains in both integrals.
\end{proof}

\begin{corollary}[Necessary angular constraints from closedness]
\label{cor:r13-angular-constant-tensor-balance}
For every smooth tangent vector field $v$ on the sphere,
\[
           \int_{\mathbb S^{N-1}}e_{\mathbb S}(v):d\mathsf M=0,
           \qquad e_{\mathbb S}(v)=\operatorname{sym}\nabla_{\mathbb S}v.
 \tag{SR14a}
\]
Equivalently $\operatorname{div}_{\mathbb S}\mathsf M=0$
in the distributional tangent sense.  The matrix also satisfies
\[
       \int_{\mathbb S^{N-1}}d\mathsf M
            =\int_{\mathbb S^{N-1}}\theta\theta^t\,d\sigma.
 \tag{SR14}
\]
Equivalently, for every $a,b\in\mathbf R^N$,
\[
 \int (a^t d\mathsf M b)
                =\int(a\cdot\theta)(b\cdot\theta)\,d\sigma.
\]
The finite moment identity alone is not asserted sufficient.
The distributional condition is used in the real converse
below; original analytic realization is a separate question.
\end{corollary}

\begin{proof}
On the punctured real space put $b(r\theta)=r v(\theta)$,
using the Euclidean identification of vectors and covectors,
and let $E=e(b)$ be its symmetric derivative.  Since
$v\perp\theta$, direct differentiation gives
$E_{rr}=0$ and $E_{\rm tan}=e_{\mathbb S}(v)$.  The tensor
$E$ is full-exact: in the retained sign convention it is
\[
 E=d'\!\left(\tfrac12\sum_jb_jd''x_j\right)
           +d''\!\left(-\tfrac12\sum_jb_jd'x_j\right).
\]
Both closedness equations of $T_0$ therefore imply
$\langle T_0,E\wedge D\rho\rangle=0$.  To justify the
integration by parts, multiply the displayed primitive
forms by a smooth cutoff equal to one on a neighborhood
of the compact annular support of $D\rho$; the added
derivatives then have zero product with $D\rho$.
Apply SR11 to obtain SR14a.  This use of SR11 is legitimate
also for a smooth degree-zero homogeneous tensor on the
puncture: its proof needs only boundedness on the sphere
and the radial formula, and $D\rho$ is supported away
from zero.

Take $E$ to be any constant symmetric tensor.  It is smooth
and closed.  Separate closedness of $T_0$, tested against
the compact form $E\wedge d''\rho$, gives
$\langle T_0,E\wedge D\rho\rangle=0$.
Apply SR11 and vary the symmetric matrix $E$.
\end{proof}

\subsection{The actual leading residual and arbitrary cutoffs}

Return to the PA--MT--LE--PT setting, with the same original
degree-one background $S$.  Let $Z_0$ be LE's full paired
leading tensor and $T_0$ any of its realized critical
background tangents.  Write
\[
 Z_0=\tfrac12\sum_\alpha c_\alpha
       \bigl(d'_P h_\alpha\wedge d''_P k_\alpha
                +d'_P k_\alpha\wedge d''_P h_\alpha\bigr),
 \tag{SR15}
\]
where $h_\alpha,k_\alpha$ are the actual ambient homogeneous
PL first jets.  Their spherical derivatives mean the
tangential derivatives of the restrictions to the sphere,
with all face restrictions retained.

\begin{theorem}[Spherical expression for the central coefficient]
\label{thm:r13-spherical-realized-residual}
With the normalization SR6,
\[
\boxed{
c_S(P)=\sum_\alpha c_\alpha
  \left(\int_{\mathbb S^{N-1}}h_\alpha k_\alpha\,d\sigma
    -\int_{\mathbb S^{N-1}}
                  d_{\mathbb S}h_\alpha\otimes
                         d_{\mathbb S}k_\alpha:d\mathsf M\right).}
 \tag{SR16}
\]
Moreover, in the retained derivative convention,
\[
                     D(Z_0\wedge T_0)=-c_S(P)\delta_0
                                  \quad\hbox{on }\mathbf R^N.
 \tag{SR17}
\]
Consequently PT's formula for $c_S(P)$ is independent of
every compact smooth cutoff equal to one near zero, whether
or not it is radial.  The angular measure may depend on the
chosen critical tangent, but its value in SR16 does not.
\end{theorem}

\begin{proof}
Euler's identity for each homogeneous scalar gives
$d h_\alpha(\theta)=h_\alpha(\theta)$ in the radial
direction.  Thus the radial entry of SR15 is
$\sum c_\alpha h_\alpha k_\alpha$, and its tangential
entry is the corresponding symmetrized spherical gradient
product.  The symmetry of $\mathsf M$ makes its contraction
with that product equal to the unsymmetrized expression
in SR16.  Apply SR11 and PT11, which supplies the negative
of that pairing.

For the stronger distributional identity, let
$\eta\in C_c^\infty(\mathbf R^N)$ be arbitrary.  For
sufficiently small $\delta>0$, proper protection makes
$\eta_\delta=\eta(z/\delta)$ an original compact smooth
test on the whole inverse image.  The actual equality
$\mathcal A_S(P)=c_S(P)\delta_y=-D M_S(Z_P)$ gives
\[
       \langle T,Z_P\wedge D[\eta(z/\delta)]\rangle
                              =-c_S(P)\eta(0).
 \tag{SR18}
\]
Here $T$ is MT's unscaled full-fiber transpose.  If the
support of $\eta$ is contained in $B_R$, LE's complete-face
Taylor estimate and MT's mass bound imply
\[
 \left|\langle T,(Z_P-Z_0)\wedge D[\eta(z/\delta)]\rangle\right|
       \le C_\eta\delta^{-1}\mu_T(B_{R\delta})
                               =O_{S,\eta}(\delta).
 \tag{SR19}
\]
Indeed the coefficient error is $O(\delta)$ there and
the Hessian is $O(\delta^{-2})$.  This uses every point
of the protected inverse image, not only the distinguished
source point.

The exact scaling from PT12 identifies
$\langle T,Z_0\wedge D[\eta(z/\delta)]\rangle$ with
$\langle T_\delta,Z_0\wedge D\eta\rangle$.
Along the subsequence defining $T_0$, PT's compatible-PS
weak continuity passes this last pairing to
$\langle T_0,Z_0\wedge D\eta\rangle$.
Equations SR18--SR19 therefore give
\[
             \langle T_0,Z_0\wedge D\eta\rangle
                                     =-c_S(P)\eta(0).
\]
This is SR17.  Taking any cutoff with value one near zero
proves its asserted independence.  No claim that the formal
quadratic first-jet presentation is itself an actual zero
has been used.
\end{proof}

\subsection{An exact real classification, without analytic realization}

\begin{theorem}[Real conic currents and divergence-free angular matrices]
\label{thm:r13-real-angular-classification}
The polar construction SR7 is a bijection between the
following two classes:
\begin{enumerate}
\item finite symmetric positive semidefinite matrix-valued
measures $\mathsf M$ on $\mathbb S^{N-1}$ satisfying
$\mathsf M\theta=0$ and
$\operatorname{div}_{\mathbb S}\mathsf M=0$;
\item symmetric middle-positive separately $d'$- and
$d''$-closed real supercurrents $T$ of bidimension $(2,2)$
on $\mathbf R^N$ with
\[
                 \lambda^2(a_\lambda)_*T=T,
                 \qquad a_\lambda(x)=x/\lambda.
 \tag{SR20}
\]
\end{enumerate}
Their traces satisfy
$\mu_T(B_R)=\tfrac12\sigma(\mathbb S^{N-1})R^2$.
This classifies real currents only.  It does not characterize
which of them are critical tangents of original allowed
non-Archimedean backgrounds under all actual refinements.
\end{theorem}

\begin{proof}
Given $\mathsf M$ in the first class, define its radial
cocoefficient measure by
\[
           d\mathsf Q=r\,dr\,A_\theta d\mathsf M A_\theta^*.
 \tag{SR21}
\]
It has locally finite trace, is positive semidefinite,
and defines a middle-positive coefficient-Radon current
by SR4.  It is symmetric, satisfies SR20, and has the
displayed trace normalization.  It remains to prove both
closedness equations.

First suppose $d\mathsf M=B(\theta)\,d\omega(\theta)$,
where $B$ is a smooth symmetric tangent matrix and $d\omega$
is Euclidean spherical surface measure.  Relative to
Lebesgue measure, SR21 has cocoefficients
\[
 C_{ij,kl}(x)=r^{2-N}
  (\theta_i\theta_kB_{jl}-\theta_i\theta_lB_{jk}
          -\theta_j\theta_kB_{il}+\theta_j\theta_lB_{ik}).
 \tag{SR22}
\]
Extend the indices antisymmetrically from $i<j$ and $k<l$.
In these complementary cocoefficients, $d'T=0$ is exactly
$\sum_i\partial_i C_{ij,kl}=0$; the second closedness
equation is its column transpose.  If
$v=\operatorname{div}_{\mathbb S}B$, direct differentiation
using $B\theta=0$ and symmetry gives
\[
 \sum_i\partial_i C_{ij,kl}
       =r^{1-N}\theta_j(\theta_l v_k-\theta_k v_l).
 \tag{SR23}
\]
For completeness,
$\partial_i\theta_j=r^{-1}(\delta_{ij}-\theta_i\theta_j)$,
and
\[
 \sum_i\partial_i B_{il}
        =r^{-1}\bigl(v_l-\theta_l\operatorname{tr}B\bigr).
\]
In differentiating SR22, the first two terms contribute
$r^{1-N}(\theta_k B_{jl}-\theta_l B_{jk})$.
The derivatives of the last two terms cancel these,
the symmetric $B_{kl}$ terms, and the trace terms,
leaving precisely SR23.  Thus $v=0$ proves both
closedness equations on the punctured space.

For a finite matrix measure, average its tensorial rotation
translates against a smooth approximate identity of $SO(N)$.
This produces smooth tangent symmetric positive semidefinite
matrix densities, converging weakly to $\mathsf M$.
Distributional spherical divergence commutes with isometries,
so every average remains divergence-free.  The smoothing
assertion follows by differentiating the group convolution;
the rotation action is smooth and transitive on the sphere.
Their currents converge coefficientwise on every compact
annulus by SR21.  Distributional continuity of $d'$ and
$d''$ therefore proves closedness of the limiting current
on the punctured space.

To include the origin, use a smooth deletion cutoff which
changes from zero to one on $\epsilon<r<2\epsilon$ and
has first derivative bounded by $C/\epsilon$.  The error
in either closedness equation, against any fixed compact
test, is bounded by
\[
            C\epsilon^{-1}\mu_T(B_{2\epsilon})
                                     =O(\epsilon).
 \tag{SR24}
\]
It tends to zero.  Thus both equations hold globally.

Conversely, for any current in the second class, SR20 gives
$\mathsf Q(\lambda E)=\lambda^2\mathsf Q(E)$.
Its trace therefore has no atom at zero and has ball mass
$\nu R^2$ for a finite $\nu$.  Apply MT's annular
Stokes identity to this positive closed real current.
Its constant normalized ball mass forces the tangential
two-vector block to vanish, and positivity removes the
mixed blocks.  This is the radial support used in SR5.
Proposition~\ref{prop:r13-angular-critical-matrix} now gives
the unique tangent positive angular measure.  Corollary
\ref{cor:r13-angular-constant-tensor-balance} gives its
distributional spherical divergence zero.  The two explicit
polar constructions are inverse.
\end{proof}

\paragraph{Scope of the angular representation.}
SR7 and SR11 are real coefficient identities for the supplied
critical tangent.  SR14a and SR14 use its closedness and supply
necessary angular equilibrium constraints.  SR16--SR17
use the actual original background, the actual zero relation,
and the complete protected-fiber PA--MT--LE--PT comparison.
They do not construct an original background from an arbitrary
positive angular measure, identify that measure with a residue
numerical class, or prove SR16 is zero.  The remaining question
is the annihilation of this full paired angular expression
for the measures realized by the original allowed backgrounds.
The real classification SR20--SR24 does not remove that
restriction or assert that every real angular measure in
its first class is analytically realized.

\paragraph{Dependencies.}
MT supplies SR1, SR4, SR8 and radial matrix support.  PT supplies
the canonical compatible-PS pairing and its weak continuity.
LE supplies the complete leading tensor and the coefficient
estimate used in SR19; PA supplies the actual central current
identity in SR18.  The polar measure construction, normalization,
and the integration in SR11 are proved in this section.
